\chapter{Spherical-Harmonic Operators}
\label{chap:sh}

The horizontal operators ($\nabla$, $\Lap$, $\Lapinv$, Helmholtz,
Jacobian) are computed via \textbf{spherical-harmonic} (SH) transforms
delegated to \texttt{pvtend.sh\_ops}.  This chapter documents the
wrappers in \texttt{sh\_ppvi.operators}, in particular how the
hemispheric ERA5 grid is extended to a full sphere by
\emph{parity mirroring}.

\section{Hemispheric Grid and Parity Mirroring}

The working grid is a Northern-Hemisphere lat-lon patch
($10.5^\circ$N--$85.5^\circ$N, $-169.5^\circ$E--$-40.5^\circ$E,
$\Delta\phi=\Delta\lambda=1.5^\circ$).  SH transforms, however,
require a \emph{full sphere}.  The function
\texttt{parity\_mirror(field, lat, kind)} extends a NH field across
the equator with prescribed symmetry:

\begin{itemize}\setlength\itemsep{0pt}
\item \texttt{kind="scalar"}: even reflection,
      $f(-\phi,\lambda)=+f(\phi,\lambda)$ —
      used for $\psi$, $\Phi$, $q$;
\item \texttt{kind="v"} (vector meridional component):
      odd reflection,
      $f(-\phi,\lambda)=-f(\phi,\lambda)$ — used for the RHS of
      divergence-like Poisson problems.
\end{itemize}

\noindent After the spectral operation the SH field is inverse-transformed
back to grid space and restricted to the NH.  This is hidden inside
each wrapper.

\section{Horizontal Operators}

All wrappers act level-by-level on $(K,n_{\text{lat}},n_{\text{lon}})$
arrays.

\paragraph{Gradient.} (\texttt{operators.gradient})
\begin{equation}
\frac{\partial f}{\partial x}=\frac{1}{a\cos\phi}\frac{\partial f}{\partial\lambda},
\qquad
\frac{\partial f}{\partial y}=\frac{1}{a}\frac{\partial f}{\partial\phi}.
\end{equation}

\paragraph{Laplacian.} (\texttt{operators.laplacian})
\begin{equation}
\Lap f = \frac{1}{a^{2}\cos^{2}\phi}\frac{\partial^{2}f}{\partial\lambda^{2}}
       + \frac{1}{a^{2}\cos\phi}\frac{\partial}{\partial\phi}\!
         \left(\cos\phi\,\frac{\partial f}{\partial\phi}\right).
\end{equation}
In spectral space, $\widehat{\Lap f}_{l,m}=-\lambda_l\,\hat f_{l,m}$
with the eigenvalues
\begin{equation}
\lambda_l=\frac{l(l+1)}{a^{2}},\qquad l=0,1,2,\ldots
\label{eq:laplacian_eigs}
\end{equation}

\paragraph{Poisson inversion.} (\texttt{operators.laplacian\_inv})
\begin{equation}
\Lap\chi=\mathrm{rhs}\;\Longrightarrow\;
\hat\chi_{l,m}=
\begin{cases}
  -\hat{\mathrm{rhs}}_{l,m}/\lambda_l, & l\geq 1,\\
  0,                                    & l=0\;\text{(mean-zero gauge)}.
\end{cases}
\label{eq:laplacian_inv}
\end{equation}

\paragraph{Helmholtz decomposition.} (\texttt{operators.helmholtz})
\begin{equation}
(u,v) = \nabla^{\perp}\psi + \nabla\chi
\;\Longleftrightarrow\;
\psi=\Lapinv\zeta,\quad \chi=\Lapinv\delta,
\end{equation}
where $\zeta=\partial v/\partial x-\partial u/\partial y$ and
$\delta=\partial u/\partial x+\partial v/\partial y$.

\paragraph{Jacobian.} (\texttt{operators.jacobian})
\begin{equation}
J(a,b)=\frac{\partial a}{\partial x}\frac{\partial b}{\partial y}
      -\frac{\partial a}{\partial y}\frac{\partial b}{\partial x}.
\end{equation}

\paragraph{$\nabla\!\cdot\!(f\nabla\psi)$.} (\texttt{operators.div\_f\_grad})
Because $f=f(\phi)$ only,
\begin{equation}
\nabla\!\cdot\!(f\nabla\psi)=f\Lap\psi+\beta\,\frac{\partial\psi}{\partial y}.
\end{equation}

\begin{keybox}[Vertical x horizontal separability]
Every horizontal SH operator acts level-by-level; every vertical
operator ($D_1$, $D_2$) acts column-by-column.  The full PPVI operator
\cref{eq:F_op} is therefore a sum of \emph{separable} actions, which
is what makes a level-mean spectral preconditioner of the vertical
operator possible.
\end{keybox}
